% ============================================================================
% Merged oral presentation — IM06 project / MedVAE
% Order of the parts: 1) Elias  2) Théophile  3) Yanic  4) Théo
% All figures are in figures/<firstname>/...
% ============================================================================
\documentclass[aspectratio=169]{beamer}

% ---------- Packages ----------
\usepackage[utf8]{inputenc}
\usepackage[T1]{fontenc}
\usepackage[english]{babel}
\usepackage{graphicx}
\usepackage{amsmath,amssymb}
\usepackage{hyperref}
\usepackage{pifont}
\usepackage{booktabs}
\usepackage{xcolor}
\usepackage{colortbl}

\newcommand{\cmark}{\textcolor{green!50!black}{\ding{51}}}
\newcommand{\xmark}{\textcolor{red!70!black}{\ding{55}}}
\DeclareMathOperator*{\argmin}{arg\,min}

% Look for figures in all subfolders
\graphicspath{%
  {figures/}%
  {figures/elias/}{figures/theophile/}{figures/yanic/}{figures/theo/}%
}

% ---------- Theme ----------
\usetheme{Madrid}
\usecolortheme{default}

\definecolor{TPBordeaux}{RGB}{140,30,45}
\setbeamercolor{structure}{fg=TPBordeaux}

% Automatic section page
\AtBeginSection[]{%
  \begin{frame}[plain,c]
    \begin{center}
      \usebeamercolor[fg]{structure}%
      {\Large\bfseries\insertsection}
    \end{center}
  \end{frame}%
}

% ---------- Title info ----------
\title[MedVAE — Coronary segmentation]{MedVAE: adaptation, fine-tuning and analysis for medical imaging and coronary segmentation}
\author{Yanic Rothlingshofer  \and Théo Palagi \and Elias Corlou \and Théophile Nadiedjoa}
\institute[]{\large \it Télécom Paris}
\date{\underline{Supervisor:} Elsa Angelini}

% ---------- Logo (shown if logo_TP.png is next to the .tex) ----------
\titlegraphic{%
  \vspace{0cm}
  
\includegraphics[width=0.3\linewidth]{logoTP.png}%
  }%


\begin{document}

\maketitle

% ############################################################################
% INTRODUCTION (from final.tex, Sec. Introduction)
% ############################################################################
\begin{frame}{Introduction}
  \begin{columns}[T]
  \begin{column}{0.58\textwidth}
  \textbf{Context.} Deep learning in medical imaging is limited by the high
  resolution required for diagnostic detail. VAEs compress the image into a
  compact latent.
  \begin{itemize}
    \item \textbf{MedVAE}: generic medical compressor ($>1$M images),
      $\times 16$ ($512^2 \to 128^2 \times 1$), trained in two stages
      (reconstruction then semantic enrichment).
  \end{itemize}
  \vspace{0.4em}
  \textbf{Domain.} Coronary angiography (ARCADE, $1\,200$ images):
  thin, tubular, low-contrast structures, noise and artifacts.
  \end{column}
  \begin{column}{0.42\textwidth}
  \begin{block}{Central question}
  Does a generic medical compressor remain \alert{reliable and useful} when
  facing a specific and demanding clinical domain?
  \end{block}
  \vspace{0.2em}
  \footnotesize \textbf{Four complementary axes:}
  \begin{enumerate}
    \item Robustness to degradations
    \item Image-quality inductive bias via FiLM
    \item Semantic structuring of the latent (JEPA)
    \item Latent representation for segmentation
  \end{enumerate}
  \end{column}
  \end{columns}
\end{frame}

% ############################################################################
% PART 1 — ELIAS
% Figures available in figures/elias/ :
%   degradation_visual.png, pipeline.pdf, mpsnr.pdf,
%   179_clean.png, 179_clean_ov.png, 179_deg_ov.png, 179_poisson.png,
%   metrics_by_level_all_masked.pdf,
%   metrics_by_level_masked_blur_NEW.png,
%   metrics_by_level_masked_jpeg_NEW.png,
%   metrics_by_level_masked_poisson_NEW.png
% ############################################################################
\section{Robustness of MedVAE to input degradations}

\begin{frame}{Method: evaluation protocol}
\centering
\includegraphics[width=0.85\linewidth]{elias/pipeline.pdf}
\end{frame}

\begin{frame}{Why not a no-reference metric?}
\begin{columns}[T]
\begin{column}{0.56\textwidth}
\small
Discarded approach: score the input quality with ARNIQA (no-reference), then correlate it with the reconstruction quality.

\vspace{1.5em}
\centering
\footnotesize
\begin{tabular}{@{}ll@{}}
\toprule
\multicolumn{2}{l}{\textbf{ARNIQA — clean dataset (n=1000)}} \\
\midrule
min / max & 0.2695 / 0.6512 \\
mean $\pm$ std & 0.4570 $\pm$ 0.0482 \\
quantiles 5/50/95 & 0.3772 / 0.4566 / 0.5391 \\
\bottomrule
\end{tabular}

\vspace{1.4em}
\raggedright
\alert{Significant but weak correlation}: the intrinsic quality variance between clean images is too small to serve as an exploitable axis.
\end{column}
\begin{column}{0.44\textwidth}
\centering
\includegraphics[width=0.68\textwidth]{elias/mpsnr.pdf}

\vspace{0.3em}
\tiny mPSNR illustration (vessel mask)
\end{column}
\end{columns}
\end{frame}

\begin{frame}{Result: three distinct regimes}
\begin{columns}[T]
\begin{column}{0.42\textwidth}
\centering
\footnotesize
\vspace{0.7em}
\begin{tabular}{lcc}
\toprule
Degradation & $r$ (Pearson) & $\rho$ (Spearman) \\
\midrule
Blur & $-0.998$ & $-0.998$ \\
Poisson & $+0.998$ & $+1.000$ \\
JPEG (before peak) & $-0.990$ & $-1.000$ \\
JPEG (after peak) & $+0.990$ & $+1.000$ \\
JPEG (global) & $(+0.030)$ & $(-0.439)$ \\
\bottomrule
\end{tabular}
\end{column}
\begin{column}{0.58\textwidth}
\footnotesize
\begin{itemize}
  \item \textbf{Blur}: near-perfect anti-correlation (28.9 $\to$ 39.2 dB) — spectral bias
  \item \textbf{Poisson}: positive correlation, mirror image — the latent removes high-frequency noise ($+1.94$ dB toward the clean image, insufficient)
  \item \alert{JPEG}: non-monotonic — smoothing that helps up to the peak, then collapse (block artifacts)
\end{itemize}
\end{column}
\end{columns}
\vspace{0.1em}
\centering
\includegraphics[width=1.0\textwidth]{elias/metrics_by_level_masked_all.pdf}
\end{frame}

% ############################################################################
% PART 2 — THÉOPHILE
% Image-quality inductive bias for MedVAE (conditioning via FiLM)
% Figures in figures/theophile/
% ############################################################################
\section{Image-quality inductive bias (FiLM)}

% Motivation
\begin{frame}{Motivation}
  \begin{itemize}
    \item MedVAE applies the \textbf{same transform} to a clean or a heavily
      degraded coronary angiography: no notion of \emph{how degraded} the
      input is.
    \item \textbf{Hypothesis:} conditioning MedVAE on a quality score
      $c \in [0,1]$, injected via FiLM, should let it adapt to the
      degradation level actually present.
    \item Three independent ways to compute $c$ (A: weighted, B: PCA,
      C: learned) $\rightarrow$ run the \emph{same} pipeline three times.
  \end{itemize}
\end{frame}

% Quality score construction
\begin{frame}{Building the quality score $c$}
  \begin{center}
    \includegraphics[width=0.95\textwidth]{theophile/pipeline_B_score.pdf}
  \end{center}
\end{frame}

% Sanity check on the score
\begin{frame}{Does the score make clinical sense?}
  \begin{center}
    \includegraphics[width=\textwidth]{theophile/result_gallery_score.png}
  \end{center}
  \begin{center}
    \footnotesize ARCADE images sorted by increasing score $c$ (approach A).
  \end{center}
\end{frame}

% FiLM conditioning + training
\begin{frame}{FiLM conditioning \& training}
  \begin{center}
    \includegraphics[width=0.95\textwidth]{theophile/pipeline_C_film.pdf}
  \end{center}
\end{frame}

% Main result
\begin{frame}{Result: conditioning hurts reconstruction}
  \begin{center}
    \includegraphics[height=0.5\textheight]{theophile/result_boxplots_eval.png}
  \end{center}
  \footnotesize
  \begin{itemize}
    \item Baseline $>$ cVAE on \textbf{all 3 approaches} $\times$
      \textbf{all 4 metrics} (MSE, SSIM, PSNR, HaarPSI), even masked
  \end{itemize}
\end{frame}

% Ablation
\begin{frame}{Ablation: does the network actually use $c$?}
  \begin{center}
    \includegraphics[width=0.85\textwidth]{theophile/result_delta_ablation.png}
  \end{center}
  \begin{center}
    \footnotesize Real $c$ vs.\ constant $c{=}0.5$ (approach C): gap grows
    with image quality: the network's response genuinely depends on $c$.
  \end{center}
\end{frame}

% ############################################################################
% PART 3 — YANIC
% Semantic structuring of the latent space (JEPA vs BioMedCLIP)
% Figures in figures/yanic/
% ############################################################################
\section{Semantic structuring of the latent space (JEPA)}

\begin{frame}{Original stage 2 — BioMedCLIP distillation}
  \centering
  \includegraphics[width=1.0\linewidth]{yanic/pipeline_ori.png}
\end{frame}

\begin{frame}{Modified stage 2 — JEPA adaptation}
  \centering
  \includegraphics[width=0.72\linewidth]{yanic/stage_2_pipeline.pdf}

  \vfill
  {\small
  \[\boxed{
    \mathcal{L}_{\text{stage2}}
    = \mathcal{L}_{\text{pred}}
    + \lambda_{\text{SIGReg}}\,\mathcal{L}_{\text{SIGReg}}
  }\]
  \[
    \mathcal{L}_{\text{pred}}
    = \frac{1}{|\mathcal{M}|} \sum_{p \in \mathcal{M}}
      \mathrm{SmoothL_1}\!\left(\hat{z}_t^{(p)},\, \tilde{z}_t^{(p)}\right),
    \qquad
    \mathcal{L}_{\text{SIGReg}}
    = \mathbb{E}_{\mathbf{v}}\!\left[ T_{\text{EP}}\!\left(\bigl\{ z_c^\top \mathbf{v} \bigr\}\right) \right]
  \]}
\end{frame}

\begin{frame}{Results — reconstruction}
  \centering
  \includegraphics[width=\linewidth]{yanic/recon.png}

  \vfill
  {\footnotesize
  \begin{itemize}
    \item \alert{Frozen} decoder: any shift of the encoder degrades the PSNR (16.85 dB for stage 2).
    \item JEPA stage 2 \alert{blurs}: predicting a masked patch amounts to \emph{averaging} the plausible regions $\Rightarrow$ loss of high frequencies (thin vessels).
    \item Consistent with the latent objective: we optimize semantics, not pixel fidelity.
  \end{itemize}}
\end{frame}

\begin{frame}{Results — RGB PCA of the intermediate features}
  \centering
  \includegraphics[height=0.82\textheight]{yanic/pca.png}
\end{frame}

% ############################################################################
% PART 4 — THÉO
% Integration of MedVAE for multi-class segmentation
% Figures in figures/theo/
% ############################################################################
\section{Integration for multi-class segmentation}

\begin{frame}{Context and problem statement}
\begin{columns}[T]
\begin{column}{0.56\textwidth}
\textbf{Task.} Fine segmentation of coronary arteries on ARCADE:
\begin{itemize}
  \item 26 classes (the background and 25 segments)
  \item thin, tubular, low-contrast structures
  \item extreme imbalance ($\approx 95\,\%$ background)
\end{itemize}
\vspace{0.6em}
\textbf{Tool.} MedVAE, a pre-trained medical variational autoencoder: $\times 16$ compression ($512^2 \to 128^2 \times 1$), but \alert{lossy}.
\end{column}
\begin{column}{0.44\textwidth}
\begin{block}{Two questions}
\begin{enumerate}
  \item Does the compression preserve enough anatomical information for segmentation?
  \item Where to integrate it: \textbf{latent} space or \textbf{pixel} space?
\end{enumerate}
\end{block}
\end{column}
\end{columns}
\end{frame}

\begin{frame}{Method: five conditions}
\begin{columns}[T]
\begin{column}{0.52\textwidth}
\small
\begin{itemize}
  \item \textbf{A} — full-resolution U-Net (reference)
  \item \textbf{A*} — U-Net from A on reconstructed images, \alert{without retraining} (control)
  \item \textbf{B} — lightweight head on the frozen latent
  \item \textbf{C} — like B, encoder fine-tuned on ARCADE
  \item \textbf{D} — U-Net retrained on reconstructed images
\end{itemize}
\vspace{0.6em}
\footnotesize \textbf{Latent} space: B, C \quad ; \quad \textbf{pixel} space: A*, D.
\end{column}
\begin{column}{0.48\textwidth}
\centering
\includegraphics[width=\linewidth]{theo/cond_B.pdf}\\[0.5em]
\includegraphics[width=\linewidth]{theo/cond_D.pdf}
\end{column}
\end{columns}
\end{frame}

\begin{frame}{Result 1: segmentation from the latent fails}
\begin{columns}[T]
\begin{column}{0.48\textwidth}
\begin{itemize}
  \item B and C: Dice $\approx 0.04$ (vs 0.433 for A)
  \item only the background is predicted, IoU $\to 0$
\end{itemize}
\textbf{Why?}
\begin{itemize}
  \item $95\,\%$ background imbalance $\Rightarrow$ trivial minimum is a strong attractor
  \item single-channel latent, \alert{no multi-scale features} (nor skip connections)
\end{itemize}
Fine-tuning the compressor (C) changes nothing.
\end{column}
\begin{column}{0.52\textwidth}
\centering
\includegraphics[width=\linewidth]{theo/comparaison_conditions.png}
\end{column}
\end{columns}
\end{frame}

\begin{frame}{Result 2: pixel space preserves, adaptation helps}
\begin{columns}[T]
\begin{column}{0.42\textwidth}
\centering
\small
\begin{tabular}{lcc}
\toprule
Cond. & Dice & IoU \\
\midrule
A   & 0.433 & 0.499 \\
A*  & 0.434 & 0.497 \\
B   & 0.038 & 0.005 \\
C   & 0.039 & 0.000 \\
D   & \textbf{0.451} & 0.495 \\
\bottomrule
\end{tabular}
\end{column}
\begin{column}{0.58\textwidth}
\small
\begin{itemize}
  \item \textbf{Compression is neutral}: A* $\approx$ A (0.434 vs 0.433)
  \item \alert{Key result: $D - A^{*} = +0.017$} \\ pure gain from adaptation (same input, only retraining differs)
  \item gain concentrated on the rare, thin segments (C11, C18, C19)
\end{itemize}
\end{column}
\end{columns}
\vspace{-0.2em}
\centering
\includegraphics[width=0.86\textwidth]{theo/dice_per_class.png}
\end{frame}

\begin{frame}{Qualitative predictions}
\vspace{0.3em}
\begin{columns}[T]
\begin{column}{0.5\textwidth}
\centering
Pixel space: structures well segmented.
\includegraphics[width=\linewidth]{theo/predAstar.png}\\[0.4em]
\includegraphics[width=\linewidth]{theo/conDmieux.png}
\end{column}
\begin{column}{0.5\textwidth}
\centering
Latent space: everything is predicted as background.
\includegraphics[width=\linewidth]{theo/predB.png}\\[0.4em]
\includegraphics[width=\linewidth]{theo/predC.png}
\end{column}
\end{columns}
\end{frame}

% ############################################################################
% GENERAL CONCLUSION (from final.tex, Sec. Conclusion)
% ############################################################################
\begin{frame}{General conclusion}
  \footnotesize
  Across four complementary axes, we probed the behavior of MedVAE,
  a generic medical compressor, on the demanding domain of coronary angiography:
  \begin{itemize}
    \item \textbf{Robustness}: intelligent low-pass filter and non-linear
      denoiser; collapses on \emph{structured} high frequencies
      (JPEG artifacts) — the nature of the degradation matters more than its intensity.
    \item \textbf{FiLM}: conditioning on a quality score slightly degrades
      the reconstruction (architectural trade-off), even though the
      network genuinely exploits the signal.
    \item \textbf{JEPA}: markedly more compact latent ($90\%$ of variance on
      $25$ dimensions), discriminative gain, but at a reconstruction cost
      (frozen decoder).
    \item \textbf{Segmentation}: $\times 16$ compression preserves
      the anatomy \alert{in pixel space} (targeted gain on thin vessels);
      the single-channel latent remains unsuited to direct dense prediction.
  \end{itemize}
  \vspace{0.3em}
  \textbf{Takeaway.} Faithful compression reusable in pixel space, but the latent
  is not yet \emph{semantically organized} nor \emph{acquisition-aware}
  enough to be segmented or conditioned directly.
\end{frame}

% ############################################################################
% PERSPECTIVES (from final.tex, Sec. Perspectives)
% ############################################################################
\begin{frame}{Perspectives}
  \begin{block}{Toward a \emph{foundation model} for radiation-based medical imaging}
  Unify the explored axes into a single training objective.
  \end{block}
  \begin{itemize}
    \item \textbf{Hybrid JEPA + reconstruction objective}: predict the masked
      latent patches while keeping the encoder anchored to a decodable space.
    \item \textbf{Physically grounded degradations}: quantum noise,
      contrast/brightness, multiplicative fields, anisotropic motion blur
      (cardiac/respiratory).
    \item \textbf{Conditioning (CVAE)} on DICOM metadata
      (modality, dose, SNR, quality) $\rightarrow$ disentangled
      content / acquisition-quality representations.
    \item \textbf{Downstream validation}: classification, detection,
      multi-modality segmentation (ARCADE as a demanding test bed).
  \end{itemize}
\end{frame}



\end{document}